\subsection{Where a second factor of seven could first occur}

\begin{claim}{A.61}{Proved}{The small character orbits cannot produce the next arithmetic factor}
Let $C/k$ be a smooth proper geometrically connected curve
of genus two over a finite field of characteristic different
from $3$ and $19$, with a rational base point.  Put
$\Delta=\pi_1(C_{\bar k})^{\{3,19\}}$ and
\[
 N_1=\Psi_3(\Delta),\qquad N_2=\Psi_{19}(N_1),
 \quad A=\Delta/N_1,\quad V=N_1/N_2,
 \quad G=\Delta/N_2,
\]
using the full mod-prime kernels of A.59.
Let $C_1,C_2$ be the associated pointed covers over $k$,
and let $\Gamma_7\simeq\mathbb Z_7$ be the canonical
pro-seven factor of $G_k$ acting through the chosen point.
Then $\Gamma_7$ acts trivially on $G$.

There is a splitting $G=V\rtimes A$, with
$A\simeq(\mathbb Z/3)^4$ and
$V\simeq(\mathbb Z/19)^{164}$.
Over $F=\mathbb Q_3(\mu_{19},\mu_3)$, the nontrivial
$V$-characters in an irreducible $G$-representation
form one $A$-orbit of length $d\in\{1,3,9,27,81\}$.
That representation has dimension $d$ and occurs in
$H^1(C_{2,\bar k},F)$ with multiplicity $2d$.
On this multiplicity space the order of the actual
$\Gamma_7$-action is at most
\[
                 7^{\,1+\lfloor\log_7(2d)\rfloor}.
                                                           \tag{A.61.1}
\]
The $V$-invariant part has trivial $\Gamma_7$-action.
Consequently
\[
 v_7\!\left(\operatorname{ord}
    (\operatorname{Frob}_k\mid H^1(C_{2,\bar k},\mathbb F_3))
          \right)\le3.                                   \tag{A.61.2}
\]
Any factor $7^2$ at this level must occur in an orbit
of length at least nine, and any factor $7^3$ must
occur in an orbit of length at least twenty-seven.
For the curve and cover of A.60, its canonical second
level therefore satisfies $1\le b_2(7)\le3$.

There is also a smaller-tower obstruction.  Start with
any pointed geometric cyclic-three cover $Y\to C$
over $k$, then take a pointed abelian nineteen-power
cover $Z\to Y$ over $k$, and finally any pointed
cover $W\to Z$ over $k$ with geometric three-group
Galois closure.  The abelian cover may have arbitrary
rank and exponent.  Then
\[
 v_7\!\left(\operatorname{ord}
    (\operatorname{Frob}_k\mid H^1(W_{\bar k},\mathbb F_3))
          \right)\le1.                                   \tag{A.61.3}
\]
Thus this extra cyclic-three step still cannot produce
the second factor of seven.  None of these bounds
asserts that a second factor actually occurs on the
full canonical tower or settles the unrestricted A12
criterion.
\end{claim}
\begin{proof}
\medskip\noindent\textbf{The geometric deck multiplicities.}
We first record precisely the deck-representation
calculation that will be used.  If $X\to C$ is a
connected finite etale Galois cover with group $H$
of order involving only $3$ and $19$, then for any
characteristic-zero splitting coefficient field $L$
of residue characteristic $3$ or $19$,
\[
 H^1(X_{\bar k},L)\simeq
                    L^{\oplus2}\oplus L[H]^{\oplus2}
                   \quad\text{as $L[H]$-modules}.          \tag{A.61.4}
\]
Only the geometric representation is compared in this
formula; no Frobenius is transported to characteristic
zero.  The proper prime-to-characteristic comparison
proved in A.4 lifts the finite cover and identifies
its geometric cohomology, with its deck action, with
that of a complex covering.  A cell structure on the
genus-two base has one vertex, four one-cells and one
two-cell.  The lifted cellular complex has modules
$L[H],L[H]^4,L[H]$ in degrees zero, one and two.
Both cohomology in degree zero and that in degree two
are the trivial representation, since the covering
is connected and the deck action preserves orientation.
The Euler identity in the representation ring gives
(A.61.4), and semisimplicity in characteristic zero
turns that identity into the displayed isomorphism.
In particular, every nontrivial absolutely irreducible
representation of degree $d$ has multiplicity $2d$;
the trivial representation has multiplicity four.

We will also use two elementary pro-prime facts,
already proved in A.59 and A.60.  For a finitely
generated pro-$r$ group, the kernel of the action of
its automorphism group on the mod-$r$ Frattini quotient
is pro-$r$.  For a free finite-rank $\mathbb Z_r$
lattice, the kernel of reduction of its automorphism
group modulo $r$ is pro-$r$.  It follows that a
pro-seven action embeds into the relevant finite
mod-$r$ image when $r=3$ or $19$.  These facts apply
to every curve and every pointed cover here, because
the full induced pro-prime topology is supplied by
the finite-core argument in A.59.

\medskip\noindent\textbf{The action on the first two deck groups is trivial.}
The mod-three Frattini quotient of the geometric
pro-three group of $C$ has rank four.  Since
$\operatorname{ord}_7(3)=6>4$, the group
$\operatorname{GL}_4(\mathbb F_3)$ has order prime to
seven.  Therefore $\Gamma_7$ acts trivially on the
entire based pro-three group of $C$.  It acts
trivially on its quotient $A$ and, by restriction,
on the pro-three group of each pointed three-group
cover of $C$.

For the full mod-three cover $C_1$, the group $A$
is elementary abelian of rank four.  The coefficient
field $\mathbb Q_{19}$ already contains $\mu_3$,
so every character of $A$ is defined over it.
Formula (A.61.4) for $C_1\to C$ decomposes its
$\mathbb Q_{19}$-cohomology into a trivial-character
space of dimension four and nontrivial-character
spaces of dimension two.  The actual $\Gamma_7$
action commutes with $A$, so it preserves these
spaces and acts linearly over $\mathbb Q_{19}$.
Its compact image preserves a lattice in each one.
Since $\operatorname{ord}_7(19)=6$, neither
$\operatorname{GL}_4(\mathbb F_{19})$ nor
$\operatorname{GL}_2(\mathbb F_{19})$ has order
divisible by seven.  Reduction is injective on a
pro-seven subgroup, so the action on all these
spaces is trivial.  It is then trivial on integral
and mod-nineteen cohomology and, by the Frattini
kernel fact, on the whole based pro-nineteen group
of $C_1$.

The arithmetic action on $G$ is consequently the
identity on both $A$ and $V$.  An automorphism of
an extension of $A$ by the abelian group $V$ that
is the identity on $A$ and on $V$ is described by
a crossed homomorphism $A\to V$: the difference
between the image of a chosen lift of $a$ and
that lift belongs to $V$, independently of the
chosen lift.  Such automorphisms form a subgroup
of the additive nineteen-group of all maps
$A\to V$.  A pro-seven group has no nontrivial
image there.  This proves that $\Gamma_7$ acts
trivially on $G$, not merely on its two factors.

The ranks follow from the full induced topology
and the unramified genus formula:
$|A|=3^4$, $g(C_1)=1+3^4=82$, and hence
$\dim_{\mathbb F_{19}}V=164$.
The extension of $A$ by $V$ splits because
$H^2(A,V)=0$: multiplication by $|A|$ is invertible
on $V$, and averaging contracts the positive-degree
finite-group cochain complex.  More explicitly,
for a normalized two-cocycle $c$, the one-cochain
$b(a)=|A|^{-1}\sum_{x\in A}c(a,x)$ satisfies
$\delta b=c$ by summing the cocycle identity.

\medskip\noindent\textbf{The multiplicity space is the only place for the new factor.}
Let $F=\mathbb Q_3(\mu_{19},\mu_3)$.
The extension obtained by adjoining $\mu_{19}$
is unramified of degree
$\operatorname{ord}_{19}(3)=18$, while adjoining
$\mu_3$ is totally ramified.  Thus the residue
field of $F$ is $\mathbb F_{3^{18}}$.

Every irreducible representation of $G=V\rtimes A$
over $F$ is described directly as follows.  Its
$V$-weights form a single $A$-orbit; otherwise the
sum of the weights in one orbit would be a proper
$G$-submodule.  Fix a weight $\lambda$ and its
stabilizer $A_\lambda$.  The character $\lambda$
extends to $V\rtimes A_\lambda$ by making it
trivial on $A_\lambda$.  Since $A_\lambda$ is
abelian and $F$ contains $\mu_3$, its irreducible
representations have degree one.  Inducing the
product of $\lambda$ with such a character gives
an irreducible representation of degree
$d=[A:A_\lambda]$; its distinct weight spaces
show irreducibility and describe every irreducible
representation.  This also proves that $F$ is a
splitting field here, with no assumption about a
Schur index.

For $\lambda\ne1$, formula (A.61.4) gives the
multiplicity $2d$.  Since $\Gamma_7$ centralizes
$G$, it acts as the identity on the fixed
irreducible factor and acts on its multiplicity
space through $\operatorname{GL}_{2d}(F)$.
The image is compact and preserves a lattice;
reduction into $\operatorname{GL}_{2d}
(\mathbb F_{3^{18}})$ is injective on this
pro-seven image.  For $j\ge1$,
\[
 v_7(3^{18j}-1)=1+v_7(j),
 \qquad v_7(3^{18}-1)=1.                                \tag{A.61.5}
\]
Indeed $\operatorname{ord}_7(3)=6$ and
$v_7(3^6-1)=1$, and the elementary geometric-series
valuation gives the displayed formula.

A finite seven-power-order matrix over
$\mathbb F_{3^{18}}$ is semisimple.  Each of its
eigenvalues lies in a field extension of degree
at most the matrix dimension.  Formula (A.61.5)
therefore bounds its order by
$7^{1+\lfloor\log_7(2d)\rfloor}$, proving
(A.61.1).  This bound uses the maximum of the
eigenvalue orders, not the seven-part of the
order of the entire general linear group.

The $V$-invariant subspace is
$H^1(C_{1,\bar k},F)$, by rational cohomological
descent, and its $\Gamma_7$-action is already
trivial from the pro-three restriction above.
The possible nontrivial orbit lengths and their
bounds are consequently
\[
\begin{array}{c|ccccc}
 d&1&3&9&27&81\\\hline
 1+\lfloor\log_7(2d)\rfloor&1&1&2&3&3.
\end{array}                                               \tag{A.61.6}
\]
The full integral cohomology of a proper smooth
curve is a free $\mathbb Z_3$-lattice, and its
mod-three cohomology is its reduction.  The
congruence kernel is pro-three, so the order of
the pro-seven image is unchanged upon this
reduction.  This proves (A.61.2) and both orbit
thresholds.  A.60 supplies the lower bound one
for its particular canonical second level.

\medskip\noindent\textbf{A cyclic-three first step remains insufficient.}
For a pointed cyclic-three cover $Y\to C$, the
same $\mathbb Q_{19}$ character argument, with
trivial multiplicity four and other multiplicities
two, proves that $\Gamma_7$ acts trivially on
the entire based pro-nineteen group of $Y$.
Given $Z\to Y$ as in the statement, take its
geometric normal core over $C$.  This core is
stable under the point-section arithmetic action
and hence gives a pointed cover over the same
field.  Its normal nineteen-group is a subgroup
of the product of the three conjugate abelian
nineteen-groups, so it is abelian.  Its quotient
is cyclic of order three.

The argument with crossed homomorphisms again
makes the $\Gamma_7$ action on this finite deck
group trivial.  Write its abelian normal group
as $V'$ and its exponent as $19^n$.  The group
splits over its cyclic-three quotient by the
same averaging argument.  Every irreducible
representation over
$F_n=\mathbb Q_3(\mu_{19^n},\mu_3)$ has degree
one or three by the weight-orbit argument.
Its nontrivial multiplicity in geometric
cohomology is therefore at most six.

We have
\[
 \operatorname{ord}_{19^n}(3)=18\cdot19^{n-1},
 \qquad v_7(3^{18\cdot19^{n-1}}-1)=1,                    \tag{A.61.7}
\]
because $v_{19}(3^{18}-1)=1$ and $7\nmid19$.
Thus the same finite-matrix bound on dimensions
at most six gives seven-part at most seven,
uniformly in $n$ and in the rank of $V'$.
The trivial $V'$-weight part is cohomology of
$Y$ and has no seven-part.  Rational pullback
and trace transfer this bound from the normal
core cover to $Z$.

Finally the Frattini automorphism kernel shows
that the pro-seven action on the full based
pro-three group of $Z$ has order at most seven.
Restriction to every pointed three-group cover,
using the full induced topology, preserves that
bound.  Its mod-three cohomology therefore
satisfies (A.61.3).
\end{proof}

\paragraph{Check and scope.}
The complete proof was independently read by the root
and finite-checks readers on 18 November 2026; both
readings passed without a requested correction.
All readers are GPT-6 Astra Pro in Pro mode, with
no separately exposed effort setting; this is
same-model-family verification and claims no priority.
It gives an upper bound and localizes the only possible
new arithmetic factor to specified multiplicity spaces.
It does not compute those spaces' actual Frobenius
matrices and does not assert that the bound is attained.
The geometric comparison A.4 is used only for the
deck-representation formula; A.59 supplies the canonical
tower and finite-core topology; A.60 supplies its
specific lower bound and the repeated elementary
Frattini argument.  No arbitrary automorphism of a
finite group is substituted for arithmetic Frobenius.
